\subsection{Parametric EPS requirements}
\label{subsection:eps-parametric-requirements}

Table~\ref{tab:eps-parameter-requirements} proposes requirements on the modeled EPS and its service to the loads. They synthesize the cited design methods using the properties in Tables~\ref{tab:eps-catalog-records}--\ref{tab:eps-specialized-properties}. Component and load limits, allocations and supplied scenarios determine the criteria. The requirements apply to the selected functions and configurations across their specified modes, temperatures, illumination, lifetime cases, schedules and pulses. Bounded calculations and supplied histories can support concept assessment through MR-03. Steady, pulse and state bounds suffice where they establish the predicates. More detailed model and software assessments can refine these results in MR-04--06.

A requirement passes when its predicate holds throughout the applicable conditions. A known violation fails. Missing decisive parameters or bounds spanning compliance and violation leave the assessment unassessed. An absent optional function is not applicable. Operating limits and absolute maxima retain their separate meanings. Pulse records determine whether startup demand replaces or supplements steady consumption. The symbol \(z_{\min,\sigma}\) denotes the minimum charge-based state of charge permitted by scenario \(\sigma\), including its applicable discharge and reserve criteria. Energy reserve criteria apply directly to stored energy.

\begingroup
\small
\setlength{\tabcolsep}{4pt}
\setlength{\parskip}{0pt}
\setlength{\LTpre}{0.4\baselineskip}
\setlength{\LTpost}{0.4\baselineskip}
\renewcommand{\arraystretch}{1.06}
\setlength{\ReportTableBodyWidth}{\dimexpr\linewidth-6\tabcolsep\relax}
\begin{longtable}{@{}P{0.09\ReportTableBodyWidth}P{0.37\ReportTableBodyWidth}P{0.29\ReportTableBodyWidth}P{0.25\ReportTableBodyWidth}@{}}
\caption{Parameter-based EPS and load requirements.}\label{tab:eps-parameter-requirements}\\
\toprule
\textbf{ID} & \textbf{Requirement and equation or algorithm} & \textbf{Pass/fail criterion} & \textbf{Rationale} \\
\midrule
\endfirsthead
\toprule
\textbf{ID} & \textbf{Requirement and equation or algorithm} & \textbf{Pass/fail criterion} & \textbf{Rationale} \\
\midrule
\endhead
\bottomrule
\endfoot
\multicolumn{4}{@{}l}{\textbf{Configuration and power paths}} \\*
EPS-01\phantomsection\label{epsreq:EPS-01} & The EPS shall use supported product configurations and settings. Compare each configured usage, including series/parallel counts, with its supported options or declared cell-assembly bounds and required external circuitry. & Every choice is supported. Configured cell counts are positive integers. Required external circuits and adapters are included. Finished panels and packs retain their fixed properties and internal arrangement. & Purchased products constrain realizable configurations \citep[Sec.~III.A--B]{review-0145}. \\
EPS-02\phantomsection\label{epsreq:EPS-02} & The solar arrays shall meet the preliminary generation bound in Equation~\eqref{eq:array-generation-screen}. Assemblies of identical cells shall also meet Equation~\eqref{eq:pv-integer-configuration}. & The sunlight generation bound holds at the attainable DET or MPPT operating point under the applicable illumination and EOL conditions. Applicable cell-count bounds hold. This screen is additional to EPS-10--12. & Illumination, operating point and EOL properties determine usable generation \citep[Secs.~3--4]{review-0289}. \\
EPS-03\phantomsection\label{epsreq:EPS-03} & The solar-array connections shall remain within the receiving input limits. Apply Equations~\eqref{eq:mppt-voltage-compatibility} and~\eqref{eq:pv-input-ratings} where applicable. For DET or shared trackers, compare the combined current--voltage characteristic at the selected operating point with the supported input range. & Operating and startup voltages remain supported. Open-circuit voltage and permissible source current satisfy their separate limits. MPPT tracking bounds and shared-input characteristics are compatible where selected. & Temperature and panel arrangements affect electrical compatibility \citep[Secs.~3--4]{review-0289}\citep[Sec.~7.1]{handpicked-0003}. \\
EPS-04\phantomsection\label{epsreq:EPS-04} & Each converter shall supply its assigned rail within its conversion envelope \(\mathcal C_r\). Compare its operating tuples with that envelope, including startup, dropout and conversion-ratio restrictions. & Every required operating tuple is supported. The converter starts at the required source condition and supplies the required output voltage and current. & Converter topology restricts attainable supply conditions \citep[Sec.~4]{handpicked-0003}. \\
EPS-05\phantomsection\label{epsreq:EPS-05} & Each charger shall support the selected battery chemistry, series count and charge policy. Compare these settings with battery and charger limits, then apply Equation~\eqref{eq:battery-charge-acceptance}. & Chemistry and cell count are supported. Charging voltage, current and termination satisfy both battery and charger rules, including charging-temperature restrictions in EPS-13. & Battery charging and regulator modes must be compatible \citep[Secs.~2.2--2.3]{review-0573}. \\
EPS-06\phantomsection\label{epsreq:EPS-06} & The EPS shall maintain each load's terminal voltage within its supply interval whenever its service is required. Apply Equations~\eqref{eq:load-path-drop}--\eqref{eq:rail-load-compatibility}, including battery-voltage bounds and pulses. Evaluate startup against its specified time and supply conditions. & Voltage limits hold throughout required service and applicable pulse intervals. Startup reaches the specified service state within the allowed time. Returns, isolation and adapters provide the required voltage relationship. & Battery and distribution drops affect load supply voltage \citep[Secs.~2.3--2.4]{review-0573}. \\
EPS-07\phantomsection\label{epsreq:EPS-07} & The EPS shall supply simultaneous steady and pulse demand within battery, converter, switch and connection current limits. Propagate power upstream with Equation~\eqref{eq:conversion-loss-accounting}, derive current at each path's voltage and sum shared-path currents. & Every input, output and shared path satisfies its continuous or pulse rating at the applicable temperature and duration. Battery charge and discharge limits remain distinct. Included consumers, losses and host execution consumption contribute once. & Coincident loads can exceed EPS current limits despite adequate average energy \citep[Sec.~VI]{review-1231}. \\
\multicolumn{4}{@{}l}{\textbf{Storage and mission energy}} \\*
EPS-08\phantomsection\label{epsreq:EPS-08} & The battery pack shall meet the preliminary capacity in Equation~\eqref{eq:battery-energy}. Derive cell-assembly ratings with Equations~\eqref{eq:battery-assembly-capacity}--\eqref{eq:battery-assembly-ratings}, or use finished-pack ratings. Use retained energy and charge in the lifetime cases. & \(E_{\mathrm{pack,BOL}}\geq E_{\mathrm{Battery}}\). Applicable EOL capacities are \(E_{\max}=f_EE_{\mathrm{pack,BOL}}\) and \(Q_{\max}=f_QQ_{\mathrm{pack,BOL}}\). Margin \(M\) enters preliminary sizing once and is not reapplied to retained capacity. Scenario sufficiency also requires EPS-10--12. & Eclipse demand and degradation constrain battery capacity \citep[Secs.~3 and~6.2.2]{handpicked-0003}. \\
EPS-09\phantomsection\label{epsreq:EPS-09} & The battery pack shall support the required nominal cycling within the selected cycle-life limits. Compare Equation~\eqref{eq:battery-mission-cycles} with the cell or pack cycle-life data at the selected DoD and operating conditions. Use Equation~\eqref{eq:battery-cycle-life} only within its applicable fit domain. & Supported cycle life is at least \(N_{\mathrm{cycle}}\). Nominal cycle DoD does not exceed \(\mathrm{DoD}_{\max}\). Exceptional survival discharge uses its separate scenario limit. & Cycle life depends on discharge depth and operating conditions \citep[Sec.~3]{handpicked-0003}. \\
EPS-10\phantomsection\label{epsreq:EPS-10} & The EPS shall supply the scheduled nominal loads while retaining their required battery reserve. Propagate discharge, accepted charging, energy and charge with Equations~\eqref{eq:battery-discharge-demand}--\eqref{eq:battery-charge-propagation}. Apply the nominal DoD condition in Equation~\eqref{eq:nominal-energy-feasibility} and scenario reserve limits. & All required demand is supplied. At every step, \(0\leq E_k\leq E_{\max}\) and \(0\leq z_{\min,\sigma}\leq z_k\leq1\). Nominal DoD and any energy reserve limits hold. Shortages fail even when orbital average generation exceeds consumption. & Energy histories expose shortages hidden by average budgets \citep[Secs.~V--VI]{review-1231}. \\
EPS-11\phantomsection\label{epsreq:EPS-11} & The EPS shall replenish the battery over each specified nominal repeating interval. Compare the propagated ending energy and charge-based state of charge with their initial values using Equation~\eqref{eq:nominal-energy-feasibility}. & \(E_{\mathrm{end}}\geq E_{\mathrm{start}}-\epsilon_E\) and \(z_{\mathrm{end}}\geq z_{\mathrm{start}}\). Here start and end identify the repeating interval boundaries. The supplied \(\epsilon_E\) is a numerical closure tolerance, not an allowance for persistent energy depletion. & Sustained operation requires generation to replenish consumed storage \citep[Sec.~3]{review-0289}. \\
EPS-12\phantomsection\label{epsreq:EPS-12} & The EPS shall supply essential loads throughout the specified survival interval and support the required recovery. Propagate the scenario's initial state, reduced illumination, load shedding and allowed transitions with the same energy and charge balances. & Essential-load service holds for the required duration. Energy and charge remain within the scenario's reserve and exceptional discharge limits. The specified recovery transition is reached within its allowed duration while satisfying EPS-04--07. & Survival following eclipse can determine battery capacity \citep[Sec.~3, Fig.~11]{review-0289}. \\
EPS-13\phantomsection\label{epsreq:EPS-13} & EPS components shall remain within their operating-temperature intervals. Compare supplied temperature histories or thermal-model bounds with each component's state-dependent limits at the corresponding cell, case or junction reference. Include required heaters as loads in EPS-07 and EPS-10--12. & Component temperature limits hold. Charging occurs only within the charging-temperature interval and discharge respects its separate interval. Required heater demand is supplied with the associated battery and load service. & Temperature restrictions and heater demand affect usable energy \citep[Sec.~6.2.3, Table~4]{handpicked-0003}. \\
EPS-14\phantomsection\label{epsreq:EPS-14} & The EPS shall dispose of available solar surplus through its supported curtailment or shunt behavior. With \(P_{\mathrm{sur},k}\) denoting unused conditioned source power (W), evaluate \(P_{\mathrm{sur},k}=\max(P_{\Delta,k},0)-P_{\mathrm{ch},k}\), including full-charge conditions. & The selected behavior accommodates the surplus within circuit voltage, current and thermal limits while maintaining required load service. Charge termination is respected. Shunted power enters dissipation. Curtailed generation is foregone harvest. & Full charge changes regulator operation and available charging power \citep[Secs.~2.2 and~3.1]{review-0573}. \\
\multicolumn{4}{@{}l}{\textbf{Protection and control states}} \\*
EPS-15\phantomsection\label{epsreq:EPS-15} & The EPS shall interrupt the declared fault currents before protected-path limits are exceeded and supply legitimate demand without nuisance trips. Compare load profiles with the no-trip region and fault current/energy/duration bounds with protected-path ratings and safe operating area. & Legitimate steady and pulse profiles remain in the no-trip region. Threshold tolerances and interruption times together keep fault stress within continuous/pulse ratings and safe operating area. Settings remain within supported protection ranges. & Protection separates legitimate demand from damaging fault stress \citep[Secs.~7.4--7.5 and~8.1.6]{handpicked-0003}. \\
EPS-16\phantomsection\label{epsreq:EPS-16} & The EPS shall enter the specified protective state at battery undervoltage or overvoltage and recover according to its hysteresis and reset policy. Apply the supplied voltage bounds and threshold crossings to the switch states, essential loads and recharge path. & Cutoff and re-enabling follow the threshold bounds, hysteresis and trip/reset delays. Output states, safe-state loads and permitted recharge paths follow the configured policy. Energized loads remain within their applicable supply limits. & Cutoff hysteresis protects the battery and permits recovery \citep[Secs.~7.4 and~8.1.6]{handpicked-0003}. \\
EPS-17\phantomsection\label{epsreq:EPS-17} & The EPS shall realize its specified power-on and inhibited states. Evaluate the applicable IF-C truth table and allowed state transitions with source-voltage bounds, startup conditions and output defaults, including specified processor/software failure states. & Each truth-table combination and allowed transition gives the required output state. Inhibited outputs remain unpowered except for explicitly permitted domains. Required autonomous protection and inhibits remain effective when the processor is unavailable. & Power-on circuitry and adopted deployment constraints define permitted energized states \citep[Sec.~7.6]{handpicked-0003}\citep[Sec.~2.3]{calpoly2022cds}. \\
EPS-18\phantomsection\label{epsreq:EPS-18} & EPS service and control ports shall support the required load-control and monitoring exchanges. Compare connected IF-D and IF-C endpoints for pin/return mapping, direction, identifiers, quantities, signal levels, protocol/version, addresses and required rate or response bounds. & Endpoint meanings and electrical levels are compatible. Addresses are unique where required by the selected transport. Required exchanges fit the supported rates and response bounds. Enable/inhibit polarity and default/fault states match the connected loads. & Telemetry and output control depend on compatible interfaces \citep[Secs.~7.2--7.5]{handpicked-0003}. \\
\multicolumn{4}{@{}l}{\textbf{Physical resource bounds}} \\*
EPS-19\phantomsection\label{epsreq:EPS-19} & The EPS shall remain within its mass and occupied-volume allocations. Apply the mass and volume inequalities in Equation~\eqref{eq:eps-product-budgets} to uniquely accounted usages, including required harnesses, adapters, packaging and protection. & Total EPS mass and occupied volume do not exceed their allocations. Aggregate and internal component values are counted once within their respective accounting boundaries. Geometric accommodation additionally requires EPS-20. & Component budgets must respect allocated resources \citep[Secs.~4--5]{review-0270}. \\
EPS-20\phantomsection\label{epsreq:EPS-20} & The EPS shall fit its allocated installation and array envelopes. Transform configured component envelopes and keep-outs into the installation frame. Compare containment, overlap, mounting, connector access and stack height for each required deployment state. Check cell-assembly area/packing and fixed-panel installation envelopes separately. & Occupied envelopes fit without unintended overlap and keep-outs remain clear. Required mounting and connector access are available. Stack height fits its allocation. For cell assemblies, active cell area does not exceed usable face area multiplied by its packing fraction. & Placement adds geometric constraints beyond volume totals \citep[Sec.~2.3]{review-0711}. \\
\end{longtable}
\endgroup
